\section{访谈建议与常见问题}
\subsection{观众提问回顾}
Graham 回答了几个经常出现的问题：Agent 能否处理非代码任务？如何衡量 Agent 成功率？什么时候需要人工接管？他的回答体现出对“人机协作”核心信念：Agent 负责重复、可编排的部分，人类保留审查与风险控制的权限。

\begin{itemize}
    \item “Agent 能否做需求分析？” → 能辅助整理文本，但仍需 PM 追问盲点。
    \item “失败时如何接管？” → Agent 生成的 diff 会被 human in the loop 复审，不通过即回滚。
    \item “如何保留 traceability？” → 所有动作都由 OpenHands 记录，日志里有时间戳、命令、结果。
\end{itemize}

\subsection{建议清单}
\begin{table}[H]
\centering
\begin{tabular}{p{0.33\textwidth} p{0.63\textwidth}}
\toprule
主题 & 建议 \\
\midrule
需求对齐 & 把需求拆成细小任务，并在 Plan 阶段写明预期输出。 \\
\midrule
日志使能 & 让 Agent 输出“意图 + 结果”到日志，便于人工分析。 \\
\midrule
安全门控 & 把数据库操作和部署命令设置为需要人工 approve 的步骤。 \\
\bottomrule
\end{tabular}
\caption{常见建议的速览。}
\end{table}

\subsection{本章小结}
\begin{itemize}
    \item 访谈 Q\&A 再次强调：Agent 不是替代，而是协助人类工程师的伙伴。
    \item 明确的建议清单便于团队把访谈内容落地。
\end{itemize}

